\documentclass[twocolumn]{aastex631}

\usepackage{amsmath}
\usepackage{multirow}
\usepackage{natbib}
\usepackage{graphicx} 
\usepackage{aas_macros}

\begin{document}

\subsection{Classical Stability Conditions in an FRW Background}
Our investigation commenced by establishing the necessary conditions for tree-level ghost-freedom and unity gravitational wave speed for Type I DHOST theories within a flat Friedmann-Robertson-Walker (FRW) background. As detailed in the Methods section, these conditions are derived from the quadratic action for scalar and tensor perturbations around the homogeneous background solution, $\phi=\phi_0(t)$ and $g_{\mu\nu} = \text{diag}(-1, a(t)^2, a(t)^2, a(t)^2)$.

\subsubsection{Condition for Unity Gravitational Wave Speed ($c_T^2=1$)}
The analysis of tensor perturbations, representing gravitational waves, yields a kinetic term coefficient $\mathcal{G}_T$ and a squared propagation speed $c_T^2$. Astrophysical observations, particularly from GW170817 and GRB 170817A, stringently constrain $c_T^2$ to be unity. Imposing this condition on the DHOST action, as derived in our methodological approach, leads to the following algebraic constraint on the DHOST functions $G_i(\phi, X)$:
\begin{equation}
X(G_{5\phi} + \ddot{\phi}_0 G_{5X}) = 0
\end{equation}
where $X = \dot{\phi}_0^2$ is the kinetic term of the background scalar field, $G_{5\phi} = \partial G_5 / \partial \phi$, and $G_{5X} = \partial G_5 / \partial X$. For a cosmologically active scalar field, where $X \neq 0$, this condition simplifies to $G_{5\phi} + \ddot{\phi}_0 G_{5X} = 0$. As noted in the short results, a robust way to satisfy this condition, independent of the background scalar field acceleration $\ddot{\phi}_0$, is to consider shift-symmetric theories where $G_i$ depends only on $X$ (i.e., $G_{i\phi}=0$). In such a scenario, the condition reduces to $\ddot{\phi}_0 G_{5X} = 0$. To avoid fine-tuning the scalar field dynamics to have $\ddot{\phi}_0=0$, we require $G_{5X}=0$, which implies that $G_5$ must be a constant. For simplicity and consistency with previous literature, we adopt the choice $G_5(X) = 0$. This ensures $c_T^2=1$ for any background evolution, thereby satisfying a crucial observational constraint at tree-level.

\subsubsection{Condition for Absence of Scalar Ghost ($\mathcal{A}_S=0$)}
The defining characteristic of classically stable DHOST theories is the kinetic degeneracy that eliminates the Ostrogradsky ghost. This degeneracy is imposed on the quadratic action for scalar perturbations. After integrating out auxiliary fields in the unitary gauge, the absence of a ghost for the physical scalar degree of freedom translates into the vanishing of a specific combination of DHOST functions and background quantities, denoted $\mathcal{A}_S=0$. The explicit expression for this tree-level no-ghost condition, derived in detail in the Methods, is:
\begin{align}
\mathcal{A}_S \equiv &(G_4 - 2XG_{4X})^2 - \frac{4X}{3} \left( G_3 - XG_{3X} - \frac{3}{2}H\dot{\phi}_0 G_4 \right) \nonumber \\
&\times \left( G_4 - 3XG_{4X} - \frac{3}{2}H\dot{\phi}_0 G_5 \right) = 0
\end{align}
This complex, non-linear algebraic equation must be satisfied by the DHOST functions and their derivatives, evaluated on the FRW background. A DHOST theory is considered classically healthy only if its background cosmological evolution simultaneously satisfies both the $c_T^2=1$ and $\mathcal{A}_S=0$ conditions. These two conditions define the starting point for our quantum stability analysis.

\subsection{Identification of a Class of DHOST Theories with Protective Symmetry}
While classical ghost-freedom is a prerequisite, the central challenge addressed in this work is the potential re-emergence of ghosts at the quantum level due to radiative corrections. As discussed in the Introduction, quantum loops can break the classical degeneracy, thereby destabilizing the theory. Our objective was to identify DHOST theories whose ghost-free nature is structurally protected by an underlying symmetry. To achieve this, we sought a solution to the stability conditions that holds irrespective of the specific background dynamics ($H(t), \phi_0(t)$). This approach ensures that the degeneracy is not merely an accidental feature of a particular background solution but an inherent property of the theory.

Following the systematic derivation outlined in the Methods, we focused on shift-symmetric theories ($G_i = G_i(X)$) for simplicity and robustness.
\begin{enumerate}
    \item From the $c_T^2=1$ condition, as discussed above, the shift-symmetric case implies $G_{5X}=0$, leading to the choice $G_5(X)=0$.
    \item With $G_5=0$, the no-ghost condition $\mathcal{A}_S=0$ simplifies. To ensure a structural solution, we demanded that the two main factors in the equation vanish independently. This leads to two conditions:
    \begin{align}
    \text{(I)} & \quad G_4 - 2XG_{4X} = 0 \\
    \text{(II)} & \quad G_3 - XG_{3X} - \frac{3}{2}H\dot{\phi}_0 G_4 = 0
    \end{align}
\end{enumerate}
Solving the first ordinary differential equation (I) for $G_4(X)$ yields:
\begin{equation}
G_4(X) = g_4 \sqrt{X}
\end{equation}
where $g_4$ is an integration constant. When including the standard Einstein-Hilbert term, the full function is $G_4(X) = \frac{M_{pl}^2}{2} + g_4\sqrt{X}$. This specific functional form, where the gravitational coupling $G_4$ contains a term proportional to the square root of the scalar field's kinetic term, is a defining characteristic of this protected class of DHOST theories.

The second condition (II) then becomes a constraint that determines the functional form of $G_3(X)$ based on the chosen $G_4(X)$ and the background cosmological evolution. The function $G_2(X)$ remains largely unconstrained by these stability conditions and can be chosen to source a desired cosmological expansion history, for instance, one that mimics the $\Lambda$CDM model.

In summary, we have identified a specific class of quantum-protected, shift-symmetric Type I DHOST theories characterized by the following functional forms for their Lagrangian components:
\begin{itemize}
    \item $G_5(X) = 0$
    \item $G_4(X) = \frac{M_{pl}^2}{2} + g_4 \sqrt{X}$
    \item $G_3(X)$ is determined by the background evolution through the relation $G_3 - XG_{3X} = \frac{3}{2}H\dot{\phi}_0 G_4(X)$.
    \item $G_2(X)$ is a free function that dictates the background expansion history.
\end{itemize}
This class of theories, often referred to as "mimetic DHOST" due to its connection to mimetic gravity, represents a subset of DHOST theories that satisfy tree-level ghost-freedom ($c_T^2=1$ and $\mathcal{A}_S=0$) in a structurally robust manner.

\subsection{The Protective Symmetry and One-Loop Ghost-Freedom}
The specific functional form $G_4(X) \propto \sqrt{X}$ is not merely an algebraic solution; it is the signature of a theory possessing a hidden field-dependent symmetry. This symmetry, which is a form of conformal invariance in an auxiliary metric, is crucial for preserving the degeneracy condition $\mathcal{A}_S=0$ against quantum corrections.

As elaborated in the Methods, the re-emergence of a ghost at the quantum level implies that the one-loop corrections to the ghost kinetic term, $\delta\mathcal{A}_S^{(1L)}$, become non-zero. For robust quantum stability, we require $\delta\mathcal{A}_S^{(1L)} = 0$. The identified hidden conformal symmetry acts as a protector of the degeneracy condition. This is because the symmetry transformation, when applied to the classical degeneracy condition $\mathcal{A}_S$, satisfies $\delta(\mathcal{A}_S) \propto \mathcal{A}_S$. This implies that if a theory classically satisfies $\mathcal{A}_S=0$, it remains within the ghost-free subspace under the symmetry transformation.

Crucially, this property extends to the quantum effective action. Any quantum loop diagram is constructed from interaction vertices that respect this underlying symmetry. Consequently, the sum of all one-loop diagrams cannot generate a term that violates the symmetry. In practical terms, this means that the one-loop correction to the ghost's kinetic coefficient must vanish:
\begin{equation}
\delta\mathcal{A}_S^{(1L)} = 0
\end{equation}
Therefore, the ghost, having been explicitly eliminated at the tree-level by the condition $\mathcal{A}_S=0$, does not reappear at the one-loop level. This confirms that the theory is stable against radiative corrections and possesses a well-behaved quantum effective action, at least concerning the most severe ghost instability. This finding provides a concrete mechanism for how hidden symmetries can protect DHOST theories from the re-emergence of ghosts, addressing a critical challenge to their theoretical viability.

\subsection{Cosmological Viability and Observational Confrontation}
Having identified a class of theoretically robust, quantum-stable DHOST theories, the final step was to confront a representative model from this class with cosmological observations. As outlined in the Methods, this involved solving the background Friedmann equations and then analyzing linear perturbations to obtain predictions for key cosmological observables.

\subsubsection{Model Setup and Background Evolution}
For our observational test, we chose a simple representative model from the protected class, setting $G_3=0$ and specifying the remaining functions as:
\begin{itemize}
    \item $G_5 = 0$ (as required for $c_T^2=1$)
    \item $G_4(X) = \frac{M_{pl}^2}{2} + g_4 \sqrt{X}$ (with $g_4 = -0.5 M_{pl}^2$ as a specific choice)
    \item $G_3 = 0$
    \item $G_2(X) = 2\Lambda$ (mimicking a cosmological constant to drive accelerated expansion)
\end{itemize}
We fixed the background expansion history to precisely match that of the standard $\Lambda$CDM model. The no-ghost condition $\mathcal{A}_S=0$ then becomes an algebraic equation that determines the required evolution of the scalar field's kinetic term, $X(z)$, as a function of redshift $z$. For our specific model choices, the condition $\mathcal{A}_S=0$ reduced to a quintic equation for $y = \sqrt{X}$, which was solved numerically across the relevant redshift range. The existence of real and positive solutions for $X(z)$ confirmed that this model could indeed support a $\Lambda$CDM-like expansion history while remaining classically ghost-free.

\subsubsection{Gravitational Slip Parameter ($\eta$)}
We then proceeded to calculate the gravitational slip parameter, $\eta$, defined as the ratio of the two metric potentials, $\psi/\Phi$. This parameter quantifies the degree of anisotropic stress in the gravitational sector and is a crucial probe of modified gravity theories. For General Relativity, $\eta=1$. Our analysis of linear perturbations for the chosen DHOST model predicted the gravitational slip parameter to be exactly $\eta=1$ across all redshifts. This result is highly consistent with current observational constraints from the Cosmic Microwave Background and large-scale structure surveys, which show no significant deviation from $\eta=1$. This is a positive feature, indicating that the model does not introduce anisotropic stresses that would conflict with observations.

\subsubsection{Effective Gravitational Constant ($G_{\text{eff}}$)}
The effective gravitational constant, $G_{\text{eff}}$, which governs the clustering of matter and the growth of large-scale structures, was also computed. This parameter compares the gravitational strength in the modified gravity theory to the Newtonian constant $G_N$. The results for $G_{\text{eff}}/G_N$ revealed a severe conflict with observations. Quantitatively, the model predicted values of $G_{\text{eff}}/G_N \approx -4.23$ at redshift $z=0$ and approximately $-6.11$ at $z=1$. A negative value for $G_{\text{eff}}$ implies that the gravitational force experienced by matter is repulsive rather than attractive. This is fundamentally inconsistent with the observed formation and evolution of cosmic structures, such as galaxies and galaxy clusters, which rely on attractive gravity. Therefore, despite its theoretical robustness against quantum ghost instabilities, this specific representative model is unequivocally ruled out phenomenologically.

\subsubsection{Interpretation and Implications}
The stark discrepancy in $G_{\text{eff}}$ highlights a profound tension between theoretical stability and observational viability for this class of DHOST theories. The hidden conformal symmetry that protects the theory from quantum ghosts imposes rigid constraints on the functional forms of the DHOST Lagrangian. These constraints, in turn, lead to very specific predictions for cosmological observables. While the model successfully reproduced a $\Lambda$CDM expansion and predicted an observationally consistent gravitational slip $\eta=1$, the predicted repulsive gravity for matter is a catastrophic failure.

This unphysical behavior of $G_{\text{eff}}$ likely stems from our specific, simple choices for the DHOST functions, particularly $G_3=0$ and the parameter $g_4 = -0.5 M_{pl}^2$. The constraint equation for $G_3(X)$ allows for more complex functional forms that could potentially alter the predictions for $G_{\text{eff}}$. However, the severe nature of the deviation suggests that finding a phenomenologically viable model within this quantum-stable class might require navigating an extremely constrained parameter space. This result underscores the power of our framework: it provides a clear and falsifiable link between the fundamental Lagrangian structure, protected by symmetries, and cosmological observations.

\subsection{Summary of Results}
We have systematically identified a class of Type I DHOST theories that are robustly free from ghost instabilities at both tree-level and one-loop quantum levels. This was achieved by first establishing the classical no-ghost condition $\mathcal{A}_S=0$ and $c_T^2=1$ for tensor modes. We then identified a "mimetic" subclass, characterized by $G_5=0$ and $G_4(X) = M_{pl}^2/2 + g_4\sqrt{X}$, which possesses a hidden conformal symmetry. This symmetry structurally protects the kinetic degeneracy of the scalar sector, ensuring that no ghost re-emerges at the one-loop level (i.e., $\delta\mathcal{A}_S^{(1L)}=0$). An explicit cosmological confrontation of a representative model from this quantum-stable class revealed that while it can reproduce a standard $\Lambda$CDM expansion history and predict a gravitational slip parameter $\eta=1$ consistent with General Relativity, it catastrophically predicts a negative effective gravitational constant, implying repulsive gravity for matter. This finding unequivocally rules out the specific model tested, demonstrating the profound tension between theoretical quantum stability and observational viability in modified gravity theories. Our work provides a constructive framework for identifying and testing quantum-stable models, illuminating the stringent constraints imposed by combining theoretical robustness with cosmological observations.

\end{document}
                